% ============================================================
% METHODOLOGY — İsmail Kabak
% CSE 476 | Deep Reinforcement Learning Based 5G Resource Allocation
% ============================================================

\section{Methodology}

\subsection{Problem Formulation as a Markov Decision Process}

We model the 5G downlink resource allocation problem as a Markov Decision
Process (MDP) defined by the tuple $(\mathcal{S}, \mathcal{A}, \mathcal{R},
\mathcal{P}, \gamma)$.

\textbf{State Space $\mathcal{S}$.}
At each time step $t$, the state vector $\mathbf{s}_t \in \mathbb{R}^{2N_{\text{UE}} + N_{\text{RB}}}$ encodes:
\begin{equation}
    \mathbf{s}_t = \bigl[\underbrace{\text{SNR}_1, \ldots, \text{SNR}_{N_{\text{UE}}}}_{\text{channel quality}},\;
                         \underbrace{T_1, \ldots, T_{N_{\text{UE}}}}_{\text{cumulative throughput}},\;
                         \underbrace{o_1, \ldots, o_{N_{\text{RB}}}}_{\text{RB occupancy}}\bigr]
\end{equation}
For the default scenario ($N_{\text{UE}}=10$, $N_{\text{RB}}=50$), the state dimension is 70.

\textbf{Action Space $\mathcal{A}$.}
A naive discrete action space of size $N_{\text{UE}}^{N_{\text{RB}}}$ is
intractable for DQN. We introduce an \emph{action adapter}: the agent selects a
priority UE $a \in \{0, 1, \ldots, N_{\text{UE}}-1\}$, and a deterministic
rule maps this integer to a full RB allocation vector of length $N_{\text{RB}}$.
This reduces the effective action dimension to $N_{\text{UE}}$.

\textbf{Reward Function $\mathcal{R}$.}
The scalar reward balances system throughput and user fairness:
\begin{equation}
    r_t = \alpha \cdot \bar{C}_t + \beta \cdot J_t
    \label{eq:reward}
\end{equation}
where $\bar{C}_t = \frac{1}{N_{\text{UE}}}\sum_{i=1}^{N_{\text{UE}}} C_{i,t}$
is the mean throughput (Mbps), $J_t$ is Jain's Fairness Index, and
$\alpha + \beta = 1$. Default: $\alpha = \beta = 0.5$.

\textbf{Transition Dynamics $\mathcal{P}$.}
State transitions are governed by the 5G channel simulation (Section~III).
The Markov property holds because the SNR state at step $t$ fully determines
the channel for the next allocation decision.

\textbf{Discount Factor $\gamma$.}
We set $\gamma = 0.99$ so the agent optimizes long-term resource utilization
rather than greedy per-step throughput.

% ------------------------------------------------------------
\subsection{Deep Q-Network Architecture}

We approximate the action-value function $Q(s, a;\,\theta)$ with a fully
connected neural network parameterized by $\theta$.

\textbf{Network structure.}
\begin{equation}
    Q(s, a;\,\theta) =
        \text{FC}_{256} \xrightarrow{\text{ReLU}}
        \text{FC}_{256} \xrightarrow{\text{ReLU}}
        \text{FC}_{N_{\text{UE}}}
\end{equation}
All linear layers are initialized with Xavier uniform initialization to
prevent vanishing and exploding gradients at the start of training.

\textbf{Target Network.}
A separate target network $Q(s,a;\,\theta^{-})$ with frozen weights is used
to compute Bellman targets. Weights are synchronized via hard update every
$C = 1000$ gradient steps:
\begin{equation}
    \theta^{-} \leftarrow \theta \quad \text{every } C \text{ steps}
\end{equation}
Without this mechanism, the optimization target shifts at every gradient step,
causing training divergence.

\textbf{Experience Replay.}
Transitions $(s_t, a_t, r_t, s_{t+1}, d_t)$ are stored in a circular buffer of
capacity $10^5$. At each training step, a mini-batch of size 64 is sampled
uniformly at random, breaking temporal correlations between consecutive
samples.

% ------------------------------------------------------------
\subsection{Training Procedure}

\textbf{Loss function.}
We minimize the mean squared Bellman error:
\begin{equation}
    \mathcal{L}(\theta) = \mathbb{E}_{(s,a,r,s',d)\sim\mathcal{B}}
        \Bigl[\bigl(y - Q(s, a;\,\theta)\bigr)^2\Bigr]
    \label{eq:loss}
\end{equation}
where the target $y$ is computed using the frozen target network:
\begin{equation}
    y = r + \gamma \cdot (1 - d) \cdot \max_{a'} Q(s', a';\,\theta^{-})
\end{equation}

\textbf{Optimizer.}
Adam optimizer with learning rate $\eta = 5 \times 10^{-4}$, selected via
grid search over $\{10^{-3},\,5 \times 10^{-4},\,10^{-4}\}$ (Section~V-A).

\textbf{Exploration schedule.}
Epsilon-greedy policy with linear decay:
\begin{equation}
    \varepsilon_t = \max\!\left(\varepsilon_{\min},\;
        \varepsilon_0 - \frac{\varepsilon_0 - \varepsilon_{\min}}{T_{\text{decay}}} \cdot t\right)
\end{equation}
where $\varepsilon_0 = 1.0$, $\varepsilon_{\min} = 0.01$, and
$T_{\text{decay}} = 50{,}000$ steps.

\textbf{Training duration.}
Each scenario is trained for 10,000 episodes with model checkpoints saved
every 100 episodes. Convergence is declared when the 100-episode rolling
average reward stabilizes (change $< 1\%$ over 500 consecutive episodes).

\textbf{Hyperparameter summary.}
\begin{table}[h]
\centering
\caption{DQN Hyperparameters}
\begin{tabular}{lll}
\hline
\textbf{Parameter} & \textbf{Value} & \textbf{Rationale} \\
\hline
Learning rate $\eta$          & $5 \times 10^{-4}$ & Best convergence in LR sweep \\
Discount factor $\gamma$      & $0.99$             & Long-horizon optimization \\
$\varepsilon$ decay steps     & $50{,}000$         & Buffer fills before decay ends \\
Target update freq $C$        & $1{,}000$          & Stable target; not too stale \\
Replay buffer size            & $100{,}000$        & Diverse sample distribution \\
Batch size                    & $64$               & Standard; balances speed/variance \\
Hidden layers                 & $2 \times 256$     & Best arch in ablation study \\
\hline
\end{tabular}
\label{tab:hyperparams}
\end{table}
